\subsection{分次模的同调维数}

在本节中，我们假设 A 是次数 ≥ 0 的分次环。我们用 \((\mathbf{A}_{n})_{n \in \mathbf{Z}}\) 表示其分次结构；因此我们有 \(A_{n} = 0\) （对 n \textless{} 0）， \(A_{0}\) 是 A 的一个子环， \(J_{0} = \bigoplus_{n > 0} A_{n}\) 是 A 的一个双边理想，并且商分次环 \(A/J_{0}\) 等同于 \(A_{0}\) .

引理 4. --- 设 M 是一个下有界的分次 A-模（X，第 56 页）。如果 \(\mathbf{A}_0\otimes_{\mathbf{A}}\mathbf{M} = 0\) ，则 \(\mathbf{M} = 0\) .

由于 \(A_{0} \otimes_{A} M\) 同构于 \(M/J_{0} M\) ，这正是 II，第 171 页，命题 6。

引理 5. --- 设 M 是一个下有界的分次 A-模，并设

\[
s: \mathbf {A} \otimes_ {\mathbf {A} _ {0}} \mathrm{M} / \mathrm{J} _ {0} \mathrm{M} \rightarrow \mathrm{M}
\]

是一个分次 A-同态，使得 \(1 \otimes_{\mathbf{A}} s: \mathbf{A}_0 \otimes_{\mathbf{A}} (\mathbf{A} \otimes_{\mathbf{A}_0} \mathbf{M} / \mathrm{J}_0 \mathbf{M}) \to \mathbf{A}_0 \otimes_{\mathbf{A}} \mathbf{M}\) 是典范同构。则 s 是满射。若 \(\operatorname{Tor}_1^{\mathbf{A}}(\mathbf{A}_0, \mathbf{M}) = 0\) ，则 s 是双射。

我们有一个正合序列

\[
0 \rightarrow \operatorname{Ker} s \rightarrow A \otimes_ {A _ {0}} M / J _ {0} M \xrightarrow {s} M \rightarrow \operatorname{Coker} s \rightarrow 0
\]

且分次 A-模 Ker \(s\) 和 Coker \(s\) 是下有界的。由正合序列 \(\mathbf{A}_{0} \otimes_{\mathbf{A}} (\mathbf{A} \otimes_{\mathbf{A}_{0}} \mathrm{M} / \mathrm{J}_{0} \mathrm{M}) \xrightarrow{1 \otimes s} \mathbf{A}_{0} \otimes_{\mathbf{A}} \mathrm{M} \to \mathbf{A}_{0} \otimes_{\mathbf{A}} \text{Coker } s \to 0\) ，我们推出 \(\mathbf{A}_{0} \otimes_{\mathbf{A}} \text{Coker } s = 0\) ，因此 \(s\) 是满射（引理 4）。于是我们有一个正合序列

\[
\operatorname{Tor} _ {1} ^ {\mathrm{A}} \left(\mathrm{A} _ {0}, \mathrm{M}\right)\rightarrow \mathrm{A} _ {0} \otimes_ {\mathrm{A}} \operatorname{Ker} s \rightarrow \mathrm{A} _ {0} \otimes_ {\mathrm{A}} \left(\mathrm{A} \otimes_ {\mathrm{A} _ {0}} \mathrm{M} / \mathrm{J} _ {0} \mathrm{M}\right) \xrightarrow {1 \otimes s} \mathrm{A} _ {0} \otimes_ {\mathrm{A}} \mathrm{M}.
\]

若 \(\operatorname{Tor}_1^{\mathbf{A}}(\mathbf{A}_0, \mathbf{M}) = 0\) ，则 \(\mathbf{A}_0 \otimes_{\mathbf{A}} \operatorname{Ker} s = 0\) 且 \(s\) 是单射（引理 4）。

命题 8. --- 设 M 是一个下有界的分次 A-模。

\begin{enumerate}
\def\labelenumi{\alph{enumi})}
\tightlist
\item
  以下条件是等价的：
\end{enumerate}

\begin{enumerate}
\def\labelenumi{(\roman{enumi})}
\item
  M 同构于如下形式的分次 A-模： \(\mathbf{A} \otimes_{\mathbf{A}_0} \mathbf{N}\) ，其中 \(N\) 是分次投射（相应地，分次自由） \(\mathbf{A}_0\) -模；
\item
  M 是投射（相应地，自由）分次 A-模；
\item
  \(\mathbf{M} / \mathbf{J}_0\mathbf{M}\) 是投射（相应地，分次自由） \(\mathbf{A}_0\) -模，且 \(\operatorname{Tor}_1^{\mathbf{A}}(\mathbf{A}_0,\mathbf{M}) = 0\) .
\end{enumerate}

\begin{enumerate}
\def\labelenumi{\alph{enumi})}
\setcounter{enumi}{1}
\tightlist
\item
  进一步假设 M 有一个由有界次数的齐次元素组成的生成系。则以下条件等价：
\end{enumerate}

\begin{enumerate}
\def\labelenumi{(\roman{enumi})}
\item
  分次 A-模 M 具有一个有限合成列，其商同构于形如 \(A \otimes_{A_{0}} N\) 的分次 A-模，其中 N 是一个平坦分次 \(A_{0}\) -模；
\item
  M 是一个平坦 A-模；
\item
  \(\mathbf{M} / \mathrm{J}_0\mathbf{M}\) 是平坦 \(\mathbf{A}_0\) -模，且 \(\operatorname{Tor}_1^{\mathbf{A}}(\mathbf{A}_0,\mathbf{M}) = 0\) .
\end{enumerate}

在这两种情形中，我们显然有 (i) \(\Rightarrow\) (ii) \(\Rightarrow\) (iii). 因此我们必须证明 (iii) \(\Rightarrow\) (i).

\begin{enumerate}
\def\labelenumi{\alph{enumi})}
\tightlist
\item
  分次 A-模 A \(\otimes_{\mathbf{A}_0}\mathbf{M} / \mathbf{J}_0\) M 是投射 A-模，因为 \(\mathbf{M} / \mathbf{J}_0\) M 是投射 \(\mathbf{A}_0\) -模。A-模的典范同态
\end{enumerate}

\[
p: \mathbf {M} \rightarrow \mathbf {M} / \mathrm{J} _ {0} \mathbf {M}
\]

是满射；因此存在一个次数为 0 的分次 A-同态

\[
s: \mathbf {A} \otimes_ {\mathbf {A} _ {0}} \mathbf {M} / \mathrm{J} _ {0} \mathbf {M} \rightarrow \mathbf {M}
\]

使得 \(p \circ s(a \otimes x) = ax\) （对 \(a \in A\) 和 \(x \in M / J_0 M\) ）.

根据引理 5， \(s\) 是从 \(\mathbf{A} \otimes_{\mathbf{A}_0} \mathrm{M} / \mathrm{J}_0 \mathrm{M}\) 到 \(\mathbf{M}\) 上的 A-模同构，由此得 (i)。

\begin{enumerate}
\def\labelenumi{\alph{enumi})}
\setcounter{enumi}{1}
\tightlist
\item
  根据对 M 的假设，存在整数 a, b 使得 \(a \leqslant b\) 且 M 由 \(\oplus M_{i}\) 生成。让我们对正整数 b - a 进行归纳推理。
\end{enumerate}

若 \(b - a = 0\) 则 M 由 \(\mathbf{M}_a\) 生成，且典范 \(\mathbf{A}_0\) -同态 \(\mathbf{M}_a \to \mathbf{M} / \mathbf{J}_0\) M 是双射的；于是由 A-同态 \(\mathbf{A} \otimes_{\mathbf{A}_0} \mathbf{M}_a \to \mathbf{M}\) （由 M 作为 A-模的结构所定义）可推出一个分次 A-同态

\[
s: \mathbf {A} \otimes_ {\mathbf {A} _ {0}} \mathbf {M} / \mathrm{J} _ {0} \mathbf {M} \rightarrow \mathbf {M}
\]

满足引理 5 的条件。于是由引理 5，s 是双射的，从而 (i) 成立。

在一般情形中，设 \(\mathbf{M}^{(a)}\) 为 M 的由 \(M_{a}\) 生成的（分次）子 A-模，且设 \(M'\) 为商 \(\mathbf{M}/\mathbf{M}^{(a)}\) 。我们有一个正合序列

\[
0 \to \mathbf {M} ^ {(a)} \xrightarrow {f} \mathbf {M} \xrightarrow {g} \mathbf {M} ^ {\prime} \to 0,
\]

由此，由于根据假设 \(\mathrm{Tor}_{1}^{\mathbf{A}}(\mathbf{A}_{0},\mathbf{M}) = 0\) ，有正合序列

\begin{enumerate}
\def\labelenumi{(\arabic{enumi})}
\setcounter{enumi}{5}
\tightlist
\item
\end{enumerate}

\[
\begin{array}{c} \operatorname{Tor} _ {2} ^ {\mathsf {A}} (\mathbf {A} _ {0}, \mathbf {M} ^ {\prime}) \to \operatorname{Tor} _ {1} ^ {\mathsf {A}} (\mathbf {A} _ {0}, \mathbf {M} ^ {(a)}) \to 0 \\ 0 \to \operatorname{Tor} _ {1} ^ {\mathsf {A}} (\mathbf {A} _ {0}, \mathbf {M} ^ {\prime}) \to \operatorname{M} ^ {(a)} / \mathrm{J} _ {0} \operatorname{M} ^ {(a)} \xrightarrow {1 \otimes f} \operatorname{M} / \mathrm{J} _ {0} \operatorname{M} \xrightarrow {1 \otimes g} \operatorname{M} ^ {\prime} / \mathrm{J} _ {0} \operatorname{M} ^ {\prime} \to 0. \end{array}\tag{7}
\]

但典范同态 \(M_{a} \to M^{(a)}/J_{0} M^{(a)}\) 是双射的。由此可知同态 \(1 \otimes f : M^{(a)}/J_{0} M^{(a)} \to M/J_{0} M\) 是单射的，并且其像是 \(A_{0}\) -子模的、 \(M/J_{0} M\) 的一个直和项。由此，从正合序列 (7) 可得 \(\operatorname{Tor}_{1}^{A}(A_{0}, M') = 0\) ，且 \(A_{0}\) -模 \(M'/J_{0} M'\) 是平坦的，因为它同构于 \(M/J_{0} M\) 的一个直和项。根据归纳假设（该假设适用于 \(M'\) ，因为后者由 \(M'^{i}\) （\(a < i \leqslant b\) ）所生成）， \(M'\) 满足条件 (i)，因此是平坦的。于是从正合序列 (6) 可推出 \(\operatorname{Tor}_{1}^{A}(A_{0}, M^{(a)}) = 0\) ，而 \(M^{(a)}/J_{0} M^{(a)}\) 等同于 \(M_{a}\) ，后者是一个平坦的 \(A_{0}\) -模（作为 \(M/J_{0} M\) 的一个直和项子模）；根据已经证明的结果，分次 A-模 \(M^{(a)}\) 同构于 \(A \otimes_{A_{0}} M_{a}\) ，因此也满足 (i)，这完成了证明。

推论 1. --- 设 M 是一个有限生成的分次 A-模。如果 \(\mathbf{A}_0\) -模 \(\mathrm{M} / \mathrm{J}_0\) M 是投射的（相应地，分次自由的，相应地，平坦的），并且如果 \(\operatorname{Tor}_1^{\mathbf{A}}(\mathbf{A}_0, \mathbf{M}) = 0\) 则 A-模 M 是投射的（相应地，分次自由的，相应地，平坦的）。

推论 2. --- 假设每个投射 \(\mathbf{A}_{0}\) -模是自由的（相应地，设 A 是诺特的，且每个有限生成的投射 \(\mathbf{A}_{0}\) -模是自由的）。设 M 是一个下有界的分次 A-模（相应地，一个有限生成的分次 A-模），并设 n 为一个整数 \(\geqslant 0\) 使得 \(\mathrm{dp}_{\mathbf{A}}(\mathbf{M}) \leqslant n\) 。存在一个由分次 A-模与次数为 0 的分次同态构成的正合序列

\[
0 \rightarrow \mathrm{L} _ {n} \rightarrow \mathrm{L} _ {n - 1} \rightarrow \dots \rightarrow \mathrm{L} _ {0} \rightarrow \mathrm{M} \rightarrow 0
\]

其中 \(L_{i}\) 是分次自由的且有下界（相应地，分次自由且有限生成）。

事实上，存在（X，第 56 页，命题 11）一个由分次 A-模与 0 次同态构成的正合序列

\[
0 \to \mathrm{L} _ {n} \to \mathrm{L} _ {n - 1} \to \dots \to \mathrm{L} _ {0} \to \mathrm{M} \to 0
\]

其中 \(L_{i}\) 对 \(0 \leqslant i \leqslant n\) 是下方有界的（相应地，有限生成的），且对 \(0 \leqslant i \leqslant n - 1\) 是分次自由的 .

由于 \(\mathrm{dp}_{\mathbf{A}}(\mathbf{M}) \leqslant n\) ，A-模 \(L_{n}\) 是投射的；因此 \(\mathbf{A}_0 \otimes_{\mathbf{A}} L_n\) 是一个投射 \(\mathbf{A}_0\) -模，从而是分次自由的；由于 \(L_{n}\) 在下方有界且 \(\mathbf{A}_0 \otimes_{\mathbf{A}} L_n\) 分次自由， \(L_{n}\) 是分次自由的（命题 7）。

推论 3（希尔伯特合冲定理）. --- 假设 \(\mathbf{A}_0\) 是一个交换域，且 A 作为 \(\mathbf{A}_0\) -代数由 \(n\) 个次数 \(>0\) 的代数无关齐次元素生成。对于每个下方有界（相应地，有限型）的分次 A-模 M，存在一个由分次 A-模和次数为 0 的分次同态构成的正合序列

\[
0 \rightarrow L _ {n} \rightarrow L _ {n - 1} \rightarrow \dots \rightarrow L _ {0} \rightarrow M \rightarrow 0,
\]

其中 \(L_{i}\) 是分次自由的且有下界（相应地，且有限生成）。

事实上，由 X，第 143 页，定理 1 有 \(\mathrm{dh}(\mathbf{A}) = n\) ，并应用推论 2。

注记. --- 推论 2 也适用于以下情形：

\begin{enumerate}
\def\labelenumi{\alph{enumi})}
\item
  \(\mathbf{A}_0\) 是主理想环且 \(\mathbf{A} = \mathbf{A}_0[\mathbf{X}_1, \dots, \mathbf{X}_{n-1}]\) ;
\item
  * A \(_{0}\) 是维数为 r 的正则局部诺特环，且 A = A \(_{0}\) {[}X \(_{1}\) , \ldots，X \(_{n}\) ，r{]}. *
\end{enumerate}

推论 4. --- 假设 \(\mathbf{A}_{0}\) 半单。设 M 为有下界的分次 A-模，n 为整数 \(\geqslant 0\) 。对于 \(\mathrm{dp}_{\mathbf{A}}(\mathbf{M}) \leqslant n\) ，其充要条件是 \(\operatorname{Tor}_{n+1}^{\mathbf{A}}(\mathbf{A}_{0}, \mathbf{M}) = 0\) .

若 \(\mathrm{dp}_{\mathbf{A}}(\mathbf{M}) \leqslant n\) ，则 \(\operatorname{Tor}_{n+1}^{\mathbf{A}}(\mathbf{A}_0, \mathbf{M}) = 0\) （X，第 135 页，引理 1）。反之，设 \(0 \to \mathbf{K} \to \mathbf{L}_{n-1} \to \ldots \to \mathbf{L}_1 \to \mathbf{L}_0 \to \mathbf{M} \to 0\) 是下方有界的分次 A-模的正合序列，使得 \(\mathbf{L}_0, \ldots, \mathbf{L}_{n-1}\) 是分次自由的（X，第 56 页，命题 11）；由 X 第 131 页的推论 3，等式 \(\operatorname{Tor}_{n+1}^{\mathbf{A}}(\mathbf{A}_0, \mathbf{M}) = 0\) 蕴含 \(\operatorname{Tor}_1^{\mathbf{A}}(\mathbf{A}_0, \mathbf{K}) = 0\) ；由于 \(\mathbf{K} / \mathbf{J}_0\mathbf{K}\) 是一个投射 \(\mathbf{A}_0\) -模（因为 \(\mathbf{A}_0\) 是半单的；X，第 140 页，命题 6），K 由命题 8（X，第 144 页）是投射的，且 \(\mathrm{dp}_{\mathbf{A}}(\mathbf{M}) \leqslant n\) .

推论 5. --- 设环 \(\mathbf{A}_0\) 半单。若 \(\operatorname{Tor}_{n+1}^{\mathbf{A}}(\mathbf{A}_0, \mathbf{A}_0) = 0\) ，则我们有 \(\mathrm{dp}_{\mathbf{A}}(\mathbf{M}) \leqslant n\) （对每个下方有界的分次 \(\mathbf{A}\) -模）。

设 A° 表示 A 的反分次环：我们有 (A°)₀ = (A₀)°，因此 (A°)₀ 是半单的（VIII，§ 5，n° 1，注记 3）。由于 \(\operatorname{Tor}_{n+1}^{A°}(A₀°, A₀°) = 0\) ，由推论 4 有 \(\mathrm{dp}_A(A₀s) \leqslant n\) ：这意味着对每个 A-模 M 有 \(\operatorname{Tor}_{n+1}^{A°}(M₀, A₀°) = 0\) ，因此 \(\operatorname{Tor}_{n+1}^{A}(A₀, M) = 0\) ，我们应用推论 4。

